\documentclass[UTF8,fontset=windows,11pt,a4paper]{ctexart}
\usepackage[margin=2.25cm]{geometry}
\usepackage{amsmath,amssymb,booktabs,tabularx,longtable,array}
\usepackage{xcolor,enumitem,fancyhdr,hyperref,xurl,float}
\setCJKmonofont{SimSun}
\definecolor{ok}{RGB}{25,120,70}
\definecolor{pending}{RGB}{190,105,10}
\definecolor{bad}{RGB}{180,35,45}
\definecolor{blue}{RGB}{35,75,130}
\definecolor{shade}{RGB}{245,247,250}
\hypersetup{colorlinks=true,linkcolor=blue,urlcolor=blue,pdfauthor={FG-MEE validation workbench},pdftitle={FG-MEE Stage 2 numerical closure report}}
\pagestyle{fancy}
\fancyhf{}
\lhead{FG-MEE 数值闭环阶段报告}
\rhead{2026-07-15}
\cfoot{\thepage}
\setlength{\headheight}{15pt}
\setlength{\parindent}{2em}
\setlength{\parskip}{0.3em}
\renewcommand{\arraystretch}{1.18}
\newcommand{\pass}{\textcolor{ok}{\textbf{完成}}}
\newcommand{\wait}{\textcolor{pending}{\textbf{待闭合}}}
\newcommand{\failure}{\textcolor{bad}{\textbf{失败}}}
\title{\textbf{FG-MEE 板双向耦合数值闭环阶段报告}\\[0.35em]\large MATLAB 三工况复算、基线漂移审计与独立 COMSOL 直驱网格验证}
\author{FG-MEE 组会验证工作台}
\date{2026年7月15日}

\begin{document}
\maketitle
\begin{center}
\fcolorbox{blue}{shade}{
\begin{minipage}{0.92\textwidth}
\textbf{阶段结论。}
当前 MATLAB 基线的 0.5、1.0、2.0 mm 反向感知结果在数值精度内严格线性；固定 300 V 与 200 A 正向探针分别得到 $-0.052257$ mm 与 $+0.174586$ mm。旧 2 mm 快照与当前绝对量差异已定位为 2026-05-22 前后的代码/案例基线漂移。新电、磁直驱模型现均可端到端运行并自动导出结果；10×10、15×15、20×20 三档网格全部完成，15×15 到 20×20 的主位移变化分别为 $0.2811\%$ 与 $0.1736\%$，通过 $0.5\%$ 网格门槛。但 20×20 结果相对 MATLAB 仍偏差 $10.7727\%$ 与 $9.2679\%$，故运行闭环已经完成，$1\%$ 交叉精度门槛尚未通过。
\end{minipage}}
\end{center}

\tableofcontents
\newpage

\section{验证目标与判据}
本阶段针对组会提出的三个核心问题推进：其一，MATLAB 当前基线能否稳定复现；其二，0.5--2.0 mm 位移下电势、磁势是否符合线性缩放；其三，300 V 与 200 A 正向工况能否由独立 COMSOL 场变量模型交叉验证。主判据为符号一致、求解正常结束、主位移相对差优先进入 $1\%$；网格变化目标为 $0.5\%$。任何求解前错误均不计为数值结果。

\section{材料与工况护照}
\begin{table}[H]
\centering
\caption{冻结的 CFFF 基准工况}
\begin{tabularx}{\textwidth}{>{\bfseries}lX}
\toprule
项目 & 冻结值 \\
\midrule
几何 & $300\,\mathrm{mm}\times300\,\mathrm{mm}\times6\,\mathrm{mm}$，10 层，每层 $0.6\,\mathrm{mm}$ \\
边界 & CFFF，$x=0$ 整侧固支 \\
材料 & $E=1.206\times10^{11}\,\mathrm{Pa}$，$\nu=0.3398$，$G=4.5\times10^{10}\,\mathrm{Pa}$ \\
压电系数 & $d_{31}=d_{32}=-5.404\times10^{-11}\,\mathrm{C/N}$ \\
压磁系数 & $q_{31}=q_{32}=49.47\,\mathrm{N/(A\,m)}$ \\
介电/磁参数 & $g_{33}=9.203\times10^{-9}\,\mathrm{F/m}$，$k_{33}=1.755\times10^{-8}\,\mathrm{H/m}$ \\
正向探针 & 外层 300 V 与 200 A，输出中心位移 \\
反向工况 & 中心位移绝对值 0.5、1.0、2.0 mm，输出电势与磁势跨度 \\
\bottomrule
\end{tabularx}
\end{table}

\section{MATLAB 三个位移工况}
\begin{table}[H]
\centering
\caption{当前基线的反向感知结果}
\small
\begin{tabular}{rrrrr}
\toprule
$|w_c|$ / mm & 电势跨度 & 磁势跨度 & 电势/mm & 磁势/mm \\
\midrule
0.5 & 164.646372 & 0.170392625 & 329.292745 & 0.340785250 \\
1.0 & 329.292745 & 0.340785250 & 329.292745 & 0.340785250 \\
2.0 & 658.585490 & 0.681570500 & 329.292745 & 0.340785250 \\
\bottomrule
\end{tabular}
\end{table}

电势的倍增残差不超过 $1.02\times10^{-12}$，磁势残差不超过 $9.99\times10^{-16}$。三次运行中的固定探针也完全一致：
\[
w_c(300\,\mathrm{V})=-0.052257\,\mathrm{mm},\qquad
w_c(200\,\mathrm{A})=+0.174586\,\mathrm{mm}.
\]
这证明当前 MATLAB 基线在所测范围内确定且线性，但并不等价于独立 COMSOL 验证完成。

\subsection{数值合理性边界}
将目标位移从 2 mm 调到 0.5 mm 后，感生电势从 658.59 降到 164.65，已避开上万伏量级；感生磁势则仅为 0.1704，2 mm 时也只有 0.6816。该磁势绝对量仍明显低于组会中引用的前人“一两百”量级，因此必须核对前人所用输出定义、参考势和工况，不能仅凭线性良好就判断物理量级正确。

\section{旧快照版本漂移}
旧 2 mm CSV 与当前文件 SHA-256 不同。旧文件生成时间早于案例文件与 \nolinkurl{matlab/run_meet_static.m} 在提交 \texttt{2f46b54} 中的变更时间。两版主要差异如下。

\begin{table}[H]
\centering
\caption{旧快照与当前 2 mm 基线比较}
\small
\begin{tabular}{lrrr}
\toprule
指标 & 旧值 & 当前值 & 当前/旧 \\
\midrule
2 mm 电驱动 / V & 90075.3299 & 11481.6945 & 0.127468 \\
2 mm 磁驱动 / A & 17974.2625 & 2291.1378 & 0.127468 \\
2 mm 感生电势 & 3882.7921 & 658.5855 & 0.169616 \\
2 mm 感生磁势 & 4.018304 & 0.681571 & 0.169616 \\
$M\!\to\!w\!\to\!E$ 归一化系数 & 0.590534615 & 0.590534615 & 1.000000000002 \\
$E\!\to\!w\!\to\!M$ 归一化系数 & 0.000121952078 & 0.000121952078 & 1.000000000010 \\
\bottomrule
\end{tabular}
\end{table}

绝对量变化具有同方向、同尺度的系统性，而归一化串联传递系数保持到数值精度，故应归类为版本基线漂移。旧快照不得直接参与当前误差统计。

\section{旧 COMSOL 模型独立性审计}
对两个 6.2 保存的 CFFF 模型进行只读容器审计，结果如下。

\begin{table}[H]
\centering
\caption{旧模型的实际物理实现}
\small
\begin{tabularx}{\textwidth}{lXX}
\toprule
模型 & 实际载荷 & 独立性判断 \\
\midrule
电模型 & 成对 $\pm4.934$ MPa 随动压力；只有 Ground，无非零 ElectricPotential & 仅等效载荷辅助复核 \\
磁模型 & 成对 $\pm16.49$ MPa 随动压力；仍为固体+静电+压电，无磁场接口 & 不能称 200 A 直接验证 \\
\bottomrule
\end{tabularx}
\end{table}

这些压力可由案例本构参数还原：
\[
dC\frac{300\,\mathrm{V}}{0.6\,\mathrm{mm}}\approx4.934\,\mathrm{MPa},\qquad
q\frac{200\,\mathrm{A}}{0.6\,\mathrm{mm}}=16.49\,\mathrm{MPa}.
\]
因此旧模型并非任意调参，但它与 MATLAB 共享同一换算链，无法排除共享理解错误。

\section{新独立模型、端到端执行与网格收敛}
新电模型直接设置两外层电极电势，并调用固体力学、静电和压电效应；新磁模型在两外层直接求解归一化磁标势拉普拉斯场，由
\[
H_z=-\frac{\partial V_m}{\partial z},\qquad
\sigma^{m}_{xx}=-q_{31}H_z,\quad \sigma^{m}_{yy}=-q_{32}H_z
\]
进入结构本构，不施加预计算面压力。两份 Java/API 驱动均已编译并完成正式求解。

300 V 电模型前两次启动后，COMSOL 日志均报告“选择 \texttt{sel\_bottom\_outer} 的几何实体层错误”，且均未生成 summary CSV。虽然系统进程返回 0，仍必须按日志判定为\failure。

第二次失败模型的内部定义证明，\texttt{sel\_bottom\_outer} 本身是合法二维边界；真正冲突来自 \texttt{sel\_outer\_electrodes}：创建 Union 时没有设置实体维度，COMSOL 将其默认成三维域选择，却接收二维边界输入。异常发生在求取 Union 计数时，尚未进入材料、物理场和网格构建，因而第一次仅修改网格引用没有生效。

经确认的第三次运行成功输出固定面 10 个边界、上下外层各 1 个域、内外电极各 2 个面，证明 Union 修正生效；随后 COMSOL 报告“不可编辑选择”。失败模型的最后历史操作与源代码顺序共同表明，问题来自试图修改默认 \texttt{lemm1} 的域选择。COMSOL 6.0 中该默认节点继承固体力学物理场选择，不可直接编辑。现已删除该操作，让外层压电材料模型通过覆盖机制接管外层，内八层继续使用默认线弹性模型。

经确认的第四次运行完成了 22,500 个扫掠六面体单元、492,969 个自由度的方程编译并开始稳态矩阵组装，随后报告默认“电荷守恒 1”缺少 \texttt{epsilonr}。这表明静电接口仍覆盖了内八个纯结构层。将 Electrostatics 限制到上下两个外压电层后，第五次运行完成 340,170 自由度线性求解，仅在恰落于外层/内层界面的电势插值点发生域归属歧义。使用已保存解将插值点向外层偏移 $10^{-9}$ m，恢复中心位移 $-0.0580492827$ mm。相同修正进入正式驱动后，第六次从头运行无错误，自动导出 $-0.0580492820$ mm。

200 A 磁模型首次正式运行即完成，15×15 网格得到中心位移 $0.1904350676$ mm、外层磁场 $H_z=\pm3.333333\times10^5$ A/m 和磁致应力 $\mp16.49$ MPa，证明磁标势、场梯度和本构应力链均按预期工作。随后完成两条链路的三档网格检查。

\begin{table}[H]
\centering
\caption{独立 COMSOL 三档网格与 MATLAB 交叉误差}
\small
\begin{tabular}{lrrrrr}
\toprule
链路 & 10×10 / mm & 15×15 / mm & 20×20 / mm & 15→20 变化 & 对 MATLAB 偏差 \\
\midrule
300 V 电致位移 & -0.0584219031 & -0.0580492820 & -0.0578865885 & 0.2811\% & 10.7727\% \\
200 A 磁致位移 &  0.1900171605 &  0.1904350676 &  0.1907662164 & 0.1736\% &  9.2679\% \\
\bottomrule
\end{tabular}
\end{table}

六个正式网格日志均无错误，六份 summary CSV 的状态均为 \texttt{completed}。两条链路已通过 $0.5\%$ 网格变化门槛，但没有通过 $1\%$ MATLAB--COMSOL 一致性门槛。继续细化网格不会解释仍稳定存在的 $9\%$--$11\%$ 差异；下一步必须审计 MATLAB 板壳广义势与 COMSOL 三维外层场/本构应力的定义口径。

\section{阶段门槛与论文表述}
\begin{table}[H]
\centering
\caption{当前验证门槛}
\begin{tabularx}{\textwidth}{lXl}
\toprule
编号 & 内容 & 状态 \\
\midrule
1 & 力致位移主链路，论文同口径主点进入 $1\%$ & \pass \\
2 & MATLAB 0.5/1/2 mm 反向链路线性与复现性 & \pass \\
3 & 300 V 电势直接 COMSOL：运行、场量与网格收敛 & \pass \\
  & 300 V 电势直接 COMSOL：对 MATLAB 进入 $1\%$ & \failure \\
4 & 200 A 磁势直接 COMSOL：运行、场量与网格收敛 & \pass \\
  & 200 A 磁势直接 COMSOL：对 MATLAB 进入 $1\%$ & \failure \\
5 & 位移到开路电势/磁势的独立 COMSOL 反向验证 & \wait \\
6 & 翘曲与功能梯度创新工况 & 基线闭环后开展 \\
\bottomrule
\end{tabularx}
\end{table}

当前可以写“MATLAB 当前基线在 0.5--2.0 mm 内严格线性；独立 COMSOL 电/磁直驱模型运行正常且达到网格收敛”；不能写“MATLAB 与 COMSOL 电/磁链误差已接近零”或“反向磁势绝对量已验证正确”。

\section{后续必要动作}
\begin{enumerate}[label=\arabic*.]
\item 逐项核对 MATLAB 板壳广义电/磁势定义与 COMSOL 三维外层场、本构应力和中心位移提取口径，解释 $9\%$--$11\%$ 系统差异；不得通过参数缩放强行追平；
\item 建立位移到开路电势/磁势的独立反向模型，冻结参考势和位移边界；
\item 对照钱沈云、赵亚飞同工况原始定义，解释磁势量级差异；
\item 精度与量级闭环后再开展翘曲和功能梯度参数研究。
\end{enumerate}

\section{可复现证据}
\begin{longtable}{>{\raggedright\arraybackslash}p{0.30\textwidth}>{\raggedright\arraybackslash}p{0.64\textwidth}}
\toprule
证据 & 路径 \\
\midrule
\endfirsthead
\toprule
证据 & 路径 \\
\midrule
\endhead
三工况 CSV/MAT/日志 & \nolinkurl{outputs/paper-20260715-fgmee/experiments/results/} \\
版本漂移对照 & \nolinkurl{outputs/paper-20260715-fgmee/experiments/current_baseline_vs_legacy.csv} \\
旧模型审计 & \nolinkurl{outputs/paper-20260715-fgmee/experiments/model_preflight_inventory.csv} \\
完整预检与命令 & \nolinkurl{outputs/paper-20260715-fgmee/experiments/stage2_preflight_audit.md} \\
电直驱代码 & \nolinkurl{tools/comsol/RunDirectElectroCfffValidation.java} \\
磁直驱代码 & \nolinkurl{tools/comsol/RunDirectMagneticCfffValidation.java} \\
三档网格汇总 & \nolinkurl{outputs/paper-20260715-fgmee/experiments/direct_comsol_mesh_validation.csv} \\
电模型最终汇总 & \nolinkurl{outputs/paper-20260715-fgmee/experiments/comsol/comsol_direct_electro_300V_20x20x10_summary.csv} \\
磁模型最终汇总 & \nolinkurl{outputs/paper-20260715-fgmee/experiments/comsol/comsol_direct_magnetic_200A_20x20x10_summary.csv} \\
正式求解日志目录 & \nolinkurl{outputs/paper-20260715-fgmee/experiments/comsol/} \\
\bottomrule
\end{longtable}

\subsection*{AI 使用说明}
本报告由 AI 辅助核查本地代码、模型容器和实验日志，并生成 LaTeX 文档。所有数值结论都绑定到列出的本地证据；独立 COMSOL 直驱结果仅证明运行与网格闭合，$1\%$ 交叉精度未通过，不作近零误差结论。
\end{document}
